%======================================================================
\section{Sky subtraction of the reduced frames}
\label{sec:cframes}

We now turn the prediction into a sky spectrum for each of the 1801 science
fibers of a reduced, flux-calibrated frame of the DRP, subtract it, and write
a sky-subtracted frame in the DRP format. A frame takes under a minute on
eight cores.

\subsection{From one prediction to every fiber}
\label{sec:perfiber}

We form the three median spectra of Section~\ref{sec:inputs} from the frame
itself, decompose them (Section~\ref{sec:decomp}), predict one vector of
science coefficients at the center of the IFU (Section~\ref{sec:ml}), and
estimate both corrections of Section~\ref{sec:postcorr} on the median
science spectrum. All fibers share these coefficients and corrections; they
differ only in their LSF, which matters, since a sky line rendered with the
nominal profile leaves a residual of about 20\% of its peak
(Section~\ref{sec:linescale}).

The sky fibers constrain the LSF best, because they contain only sky. We
decompose each Gaia-clean sky fiber (49 to 117 of the 119 in the frames
below) with its telescope's median LSF held fixed, and fit one surface per
spectrograph and channel to all of them,
\begin{equation}
  \phi(\Delta\lambda;\lambda,f)=\sum_{k,q,s}H_{kqs}\,M_k(\Delta\lambda)\,
  B_q(\lambda)\,S_s(f),
\end{equation}
with the M-splines $M_k$ of Section~\ref{sec:lsf} and four cubic B-splines
each in wavelength, $B_q$, and in the fiber position along the slit, $S_s$.
We require $H\ge0$ and $\sum_k H_{kqs}=1$, which keep the profile
nonnegative and of unit integral, and solve one quadratic program per block,
weighted by the observed inverse variance. We fit the $r$ and $z$ channels
and keep the survey LSF in $b$. No science spectrum enters the fit, so no
science signal can shape the profile used to subtract its sky. The rendered
sky is linear in the LSF coefficients, so all science fibers are rendered
by one matrix product.

The sky-subtracted frame keeps the DRP extensions, with \texttt{FLUX} the
sky-subtracted spectra and \texttt{SKY} the subtracted sky, and stores the
medians, their decompositions, the predicted coefficients, the corrections
and each fiber's LSF, from which any fiber's sky can be reconstructed. The
sky fibers are rendered with their telescope's median model; their residual
is no independent test, since they built that median and the LSF.

\subsection{The sky-arm correction over the full band}
\label{sec:fullband}

The reduced frames show a red counterpart of the solar pattern of
Section~\ref{sec:armcorr}. With the Moon up, the O$_2$ A and B bands leave
a residual of the same fractional depth in every decomposition. On a
full-Moon exposure it is $-2.4$, $-2.3$ and $-2.0\%$ of the sky continuum at
the head of the B band in the science, east and west medians, and $-9$ to
$-11\%$ at the A band, while the continuum around them is fitted to better
than 1\%; with the Moon down it is absent. We attribute it to the path of the
moonlight. The DRP divides every fiber by the transmission along the
science line of sight (Section~\ref{sec:inputs}), and the model's moonlight
has no absorption of its own, but scattered moonlight crosses more O$_2$,
along a path set mostly by the altitude of the Moon. The sky arms measure
the deficit, so we apply the correction over the whole band. On the held-out
observations, followed by the line scaling, this lowers the median $\chi^2$
from 0.139 to 0.077 at 5000--5800\,\AA, from 0.244 to 0.111 in the red and
from 0.366 to 0.125 in $z$, and on the OH-line pixels of the red from 0.493
to 0.064; 91\% of the observations improve. The fitted OH band scales move
by a median of 0.005, so the two corrections do not compete. Leaving the OH
lines out of the transfer gains nothing in the continuum and is four times
worse on the lines.

\subsection{Protecting the science lines}
\label{sec:lineprotect}

The safeguard of Section~\ref{sec:armcorr} rejects emission seen by one sky
field, not by both. Diffuse Galactic H$\alpha$, [N\,\textsc{ii}],
[S\,\textsc{ii}], [S\,\textsc{iii}], [O\,\textsc{ii}] or [O\,\textsc{iii}]
emission common to both fields, injected into both arms at the local
continuum level, reaches the science sky at 94--116\% of its flux, below
5000\,\AA{} as over the full band, and would be subtracted from the target.
We therefore treat 59 nebular lines from H12 to P8 in windows of
$\pm\max(500\,{\rm km\,s^{-1}},\,3\,{\rm FWHM})$ redward of 5800\,\AA{} and
$\pm\max(300\,{\rm km\,s^{-1}},\,2\,{\rm FWHM})$ blueward, 15\% of the
pixels. Dropping the correction there would leave the sky under the lines
uncorrected, so we keep what does not come from the lines. In the blue, where
the correction is the slowly varying solar pattern, we interpolate it across
each window between the medians of the clean pixels beside it; narrower
windows let a line at $\pm150$\,km\,s$^{-1}$ reach those pixels. In the red
the windows hold OH and [O\,\textsc{i}] airglow, which interpolation cannot
follow. There we keep the correction in proportion to the OH fraction of the
predicted sky and set it to zero in a core of $\pm(150\,{\rm
km\,s^{-1}}+1.5\,{\rm FWHM})$, where a nebular line on an OH line would
pass. For pointings in the LMC and SMC we add cores at their systemic
velocities, 262 and 146\,km\,s$^{-1}$.

Of 18 lines injected into both arms of 400 held-out observations, at rest,
at $\pm150$\,km\,s$^{-1}$ and at 262\,km\,s$^{-1}$ with the LMC cores, none
is transferred by more than 0.2\% at the 90th percentile, including
[O\,\textsc{i}]\,6300 on its airglow line and [S\,\textsc{iii}]\,9069, 9531
on OH lines. Inside the windows, on observations without nebular emission,
the median $\chi^2$ in the blue, green, red and $z$ is 0.114, 0.070, 0.204
and 0.245, against 0.240, 0.149, 0.235 and 0.309 with the correction dropped
and 0.076, 0.064, 0.095 and 0.118 without protection. We cannot center the
cores on the emission of each frame: the OH residual under H$\alpha$ in the
sky arms exceeded five times the photon noise on all 400 observations, so no
threshold detects nebular emission there. Gas beyond $\pm150$\,km\,s$^{-1}$
outside the Magellanic Clouds can therefore pass on OH lines in the red, a
quarter of its flux at the median at 260\,km\,s$^{-1}$.

\subsection{Results}
\label{sec:cframe-results}

Figure~\ref{fig:cframe} compares, for one exposure with the Moon up, the
sky subtraction of version 1.3.2 of the DRP with ours. We measure both on the
median science spectrum, outside the science-line windows, in percent of the
local sky. After our subtraction the stars of the field remain as horizontal
bands and its nebular emission as vertical lines, and faint sky residuals
remain only among the strongest OH lines beyond 9300\,\AA.

The DRP over-subtracts the continuum. Its median science spectrum lies 5.0
to 6.4\% of the sky below zero at all wavelengths, which no field can
produce. Ours lies between $+0.3\%$ at 3600--4300\,\AA{} and $+4.6\%$ in
the A band, rising to the red, as the starlight of a field at
$b=-13^\circ$ should. In the blue the solar pattern
(Section~\ref{sec:armcorr}) is removed: the RMS after subtracting a
15\,\AA{} running median falls from 1.5 to 1.0\% at 3600--4300\,\AA{} and
from 1.1 to 0.7\% at 4300--5000\,\AA. Toward the red the gain grows, to
5.8 against 1.5\% at 7720--8900\,\AA{} and 4.3 against 2.4\% in the A band.
The sky lines are removed more completely. On the sky-line pixels the
absolute residual of the median is 3.1, 3.4 and 2.6\% of the line flux in
the DRP frame at 5800--7580, 7580--8900 and 8900--9800\,\AA, and 1.5, 0.8
and 1.2\% in ours; in single fibers beyond 7580\,\AA{} it is 5.3 against
3.4\%. The DRP leaves $-4.9\%$ of [O\,\textsc{i}]\,5577 and removes
[O\,\textsc{i}]\,6300 and 6364 twice, leaving $-105$ and $-110\%$ of their
flux; we leave $-0.6$, $-6$ and $-10\%$.

In three frames
(Table~\ref{tab:cframes}) the full-band correction lowers the residual of
the sky by factors of 1.3 to 5 in the red and $z$, most at the O$_2$ bands
and with the Moon bright. It worsens the A band in one frame, whose west sky
median holds a single Gaia-clean fiber. Two residuals remain on the
full-Moon frame. The median continuum after subtraction falls from $-0.3$ to
$-1.1\%$ of the sky at 6960--7580\,\AA, and since the median of the science
fibers includes the field's own light, the sky is over-subtracted by at
least that much. The residual continuum also differs by up to 2.4\% of the
sky between the three spectrographs, each of which sees one third of the
IFU, while the sky fibers of every spectrograph agree with their medians to
within 0.8\%. A gradient of the moonlight across the half-degree field,
which one prediction at its center cannot follow, would produce such a
pattern, but we have not established the cause.

\begin{deluxetable}{lccc}
\tablecaption{High-frequency residual of the median science spectrum, in
percent of the sky continuum, with the sky-arm correction below 5000\,\AA{}
/ over the full band.\label{tab:cframes}}
\tablehead{\colhead{} & \colhead{12619} & \colhead{13020} &
\colhead{14168}}
\startdata
$h_{\rm moon}$, illumination & $4^\circ$, 0.44 & $46^\circ$, 1.00 &
$-55^\circ$, 0.11 \\
5800--6850\,\AA & 3.4 / 2.4 & 2.3 / 0.5 & 3.8 / 2.3 \\
O$_2$ B band & 8.9 / 4.2 & 2.4 / 0.6 & 11.5 / 4.0 \\
6960--7580\,\AA & 11.5 / 5.0 & 3.1 / 0.9 & 16.3 / 4.9 \\
O$_2$ A band & 6.5 / 9.4 & 3.8 / 1.5 & 8.7 / 5.1 \\
7720--8900\,\AA & 9.0 / 6.8 & 2.5 / 1.5 & 10.5 / 6.5 \\
\enddata
\tablecomments{RMS after subtracting a 15\,\AA{} running median, outside the
science-line windows; B band 6855--6960\,\AA, A band 7585--7720\,\AA.}
\end{deluxetable}

\begin{figure*}
\centering
\includegraphics[width=\textwidth]{figures/cframe_before_after_12781.pdf}
\caption{The science fibers of exposure 12781 (Moon at $31^\circ$ altitude,
illumination 0.74, $74^\circ$ from the field) before sky subtraction (top),
after the sky subtraction of DRP version 1.3.2 (middle) and after ours
(bottom). In each block, the first row shows every second fiber in four
400\,\AA{} windows, at [O\,\textsc{ii}]\,3727, H$\alpha$, the O$_2$ A band
and the OH lines around [S\,\textsc{iii}]\,9069 (marked by bars above the
second row); the second row the 200 fibers at the center of the slit over
the full range; the third the median science spectrum. Fibers are in order
along the slit, and horizontal lines separate the three spectrographs; gray
marks pixels and fibers masked in any of the three frames. The top block
uses a logarithmic scale; the two sky-subtracted blocks share one
zero-centered asinh scale, and their median spectra one asinh axis. Flux in
$10^{-14}$\,erg\,s$^{-1}$\,cm$^{-2}$\,\AA$^{-1}$. The DRP frame is
negative across the continuum, and its sky lines leave residuals, most
clearly at [O\,\textsc{i}]\,6300 and in $z$; ours leaves the stars of the
field as horizontal bands and its nebular emission as vertical lines.}
\label{fig:cframe}
\end{figure*}
